Automated resume screening is widespread and increasingly regulated, yet published matchers usually return an
opaque score, depend on private application logs, and leave bias unmeasured. We present EviMatch, an open and
deterministic pipeline that removes identity cues before scoring, parses a posting into weighted requirements,
links each requirement to the resume sentence that best supports it, and combines lexical, semantic, skill and
experience features in a low-capacity head. On a public benchmark of 8,000 labelled pairs we first show that the
official split is leaky: 99.9\% of test resumes also occur in training, and a predictor that ignores the posting
reaches macro-F1 0.42. Under a doubly-disjoint protocol, in which neither resume nor posting of a test pair is
seen in training, EviMatch reaches macro-F1 0.473 against 0.405--0.412 for TF-IDF, BM25, dense-similarity and
supervised embedding baselines (paired differences 0.06--0.07, intervals excluding zero). Ablations indicate that
experience and skill features carry the signal, whereas semantic coverage supplies evidence rather than accuracy.
Deleting the sentences an explanation cites changes the score about four times more than deleting random text.
A counterfactual audit found and led us to fix a defect in our own redactor; afterwards, name cues change the
score by exactly zero, while proxy cues such as affinity organisations still shift it by up to 0.2 points on a
100-point scale. Absolute accuracy is modest and labels are noisy, so the system supports review and should not
automate decisions.
